% Epistemic Types for Scientific Computing
% Target: PLDI 2027
%
% Compile: pdflatex main && bibtex main && pdflatex main && pdflatex main

\documentclass[acmsmall,screen,review,anonymous]{acmart}

% Packages
\usepackage{mathpartir}
\usepackage{stmaryrd}
\usepackage{booktabs}
\usepackage{subcaption}
\usepackage{xspace}

% Macros
\newcommand{\sounio}{\textsc{Sounio}\xspace}
\newcommand{\know}[1]{\texttt{Knowledge}\langle #1 \rangle}
\newcommand{\knowfull}[4]{\texttt{Knowledge}\langle #1, #2, #3, #4 \rangle}
\newcommand{\cknow}[3]{\texttt{CausalKnowledge}\langle #1, #2, #3 \rangle}
\newcommand{\eps}{\epsilon}

% Code listing style
\lstdefinelanguage{sounio}{
  keywords={fn, let, var, if, else, while, for, return, struct, enum, match,
            with, kernel, linear, type, use, effect, handle, perform},
  keywordstyle=\color{blue!60!black}\bfseries,
  ndkeywords={Knowledge, CausalKnowledge, CausalGraph, f64, i32, bool,
              IO, Mut, Alloc, GPU, Prob},
  ndkeywordstyle=\color{purple!60!black},
  sensitive=true,
  comment=[l]{//},
  commentstyle=\color{green!40!black}\itshape,
  stringstyle=\color{red!60!black},
  morestring=[b]",
}

\lstset{
  language=sounio,
  basicstyle=\ttfamily\small,
  frame=single,
  numbers=left,
  numberstyle=\tiny\color{gray},
  breaklines=true,
  captionpos=b,
  xleftmargin=2em,
  aboveskip=0.8em,
  belowskip=0.5em,
}

% ACM metadata
\setcopyright{acmcopyright}
\acmJournal{PACMPL}
\acmVolume{1}
\acmNumber{PLDI}
\acmArticle{1}
\acmYear{2027}
\acmMonth{6}
\acmDOI{10.1145/nnnnnnn.nnnnnnn}

\begin{document}

\title{Epistemic Types for Scientific Computing}
\subtitle{Type-Level Uncertainty Quantification with GUM Compliance}

\author{Anonymous}
\affiliation{%
  \institution{Anonymous Institution}
  \country{Anonymous}
}
\email{anonymous@example.org}

\begin{abstract}
Scientific computing relies on measured quantities with inherent uncertainty,
yet mainstream programming languages treat numerical values as exact.
This forces developers to either ignore uncertainty (compromising scientific
validity), track it manually (introducing propagation errors), or use library
wrappers (incurring runtime overhead without compile-time guarantees).

We introduce \emph{epistemic types}, a type system extension that makes
measurement uncertainty a first-class type-level concern.
The $\know{T}$ type encapsulates a value with its GUM-compliant standard
uncertainty, provenance chain, and temporal validity.
Arithmetic on epistemic types automatically propagates uncertainty following
ISO/IEC Guide 98-3 (the GUM), and the type system statically rejects
programs that silently discard uncertainty information.

We prove type safety (progress and preservation) and soundness of uncertainty
propagation with respect to the GUM's law of propagation of uncertainty.
We implement epistemic types in \sounio, a systems programming language,
and evaluate on three scientific computing benchmarks:
quaternionic neural networks with Bayesian uncertainty (QNN-MNIST),
physiologically-based pharmacokinetic modeling with causal interventions (PBPK),
and functional neuroimaging connectivity analysis (fMRI).
\sounio achieves $2\times$ speedup over Julia's \texttt{Measurements.jl}
and $10$--$40\times$ over Python's \texttt{uncertainties}, while providing
static guarantees that neither library can offer.
\end{abstract}

\begin{CCSXML}
<ccs2012>
<concept><concept_id>10011007.10011006.10011008.10011009.10011012</concept_id>
<concept_desc>Software and its engineering~Functional languages</concept_desc>
<concept_significance>500</concept_significance></concept>
<concept><concept_id>10011007.10011006.10011008.10011009.10011015</concept_id>
<concept_desc>Software and its engineering~Constraint and logic languages</concept_desc>
<concept_significance>300</concept_significance></concept>
</ccs2012>
\end{CCSXML}

\ccsdesc[500]{Software and its engineering~Functional languages}
\ccsdesc[300]{Software and its engineering~Constraint and logic languages}

\keywords{uncertainty quantification, type systems, scientific computing, GUM,
epistemic types, measurement uncertainty}

\maketitle

% ============================================================================
\section{Introduction}
\label{sec:intro}
% ============================================================================

Every measured quantity carries uncertainty.
A laboratory balance reports mass as $m = 10.023 \pm 0.002\,\text{g}$;
a clinical assay yields serum creatinine as $\text{SCr} = 1.2 \pm 0.1\,\text{mg/dL}$;
a brain scanner produces voxel intensities with noise proportional
to $1/\sqrt{\text{SNR}}$.
When these quantities enter computation---drug dosing calculations,
climate projections, neural network training---their uncertainties must
propagate correctly through every arithmetic operation,
or the results are scientifically meaningless.

The Guide to the Expression of Uncertainty in Measurement
(GUM)~\cite{jcgm2008gum} provides the international standard for
uncertainty propagation.
For a function $y = f(x_1, \ldots, x_N)$ of uncorrelated inputs,
the GUM specifies the combined standard uncertainty:
\begin{equation}
  u_c^2(y) = \sum_{i=1}^{N} \left(\frac{\partial f}{\partial x_i}\right)^2 u^2(x_i)
  \label{eq:gum}
\end{equation}
This first-order Taylor expansion is exact for linear functions and
provides the standard approximation for nonlinear cases.
Regulatory frameworks---FDA 21 CFR Part 11~\cite{fda21cfr11} for pharmaceutical
software, ISO 17025~\cite{iso17025} for calibration laboratories---require
GUM-compliant uncertainty reporting with auditable provenance.

Despite this critical need, mainstream programming languages offer
no native support for uncertainty.
Practitioners face three unsatisfactory options:

\paragraph{Ignore uncertainty.}
Write \texttt{let dose = weight * 0.5} and produce a point estimate
with no indication of reliability.
This is the overwhelmingly common choice, contributing to the
reproducibility crisis across sciences~\cite{eklund2016cluster}.

\paragraph{Track uncertainty manually.}
Maintain parallel variables (\texttt{dose\_mean}, \texttt{dose\_var})
and propagate by hand.
This is error-prone: a single forgotten update produces a silently
wrong uncertainty, and manual implementations frequently violate
GUM rules~\cite{jcgm2008supplement1}.

\paragraph{Use a library.}
Python's \texttt{uncertainties}~\cite{lebigot2024uncertainties} and
Julia's \texttt{Measurements.jl}~\cite{giordano2016measurements} provide
operator-overloaded uncertainty types.
However, libraries have three fundamental limitations:
(1)~no static guarantee that uncertainty is preserved---a function
accepting \texttt{float64} can silently discard it;
(2)~runtime overhead from wrapper objects, heap allocation, and
dynamic dispatch (10--40$\times$ for Python);
(3)~no provenance tracking for regulatory audit trails.

\paragraph{Our approach: epistemic types.}
We introduce \emph{epistemic types} that make uncertainty a
\emph{first-class type system construct}.
The $\know{T}$ type carries:
\begin{itemize}
  \item A point estimate $v : T$
  \item Standard uncertainty $\epsilon \in \mathbb{R}_{\geq 0}$ (GUM $u$)
  \item Provenance $\Phi$ for regulatory traceability
  \item Temporal validity $t$ (expiration for perishable measurements)
\end{itemize}

Arithmetic on $\know{T}$ automatically propagates uncertainty
using GUM Equation~\eqref{eq:gum}, and the type system \emph{statically
rejects} programs that discard uncertainty without explicit acknowledgment.
Extraction of a raw value requires compile-time proof (via SMT) or
runtime assertion that the uncertainty is below a specified threshold.

\paragraph{Contributions.}
\begin{enumerate}
  \item A type system for epistemic values with formal typing rules,
        operational semantics, and proofs of progress, preservation,
        and GUM compliance soundness (\S\ref{sec:type-system}).
  \item An implementation in \sounio, a systems programming language,
        with a Scientific IR (SIR) that optimizes epistemic computations
        via variance fusion, exact elision, and SIMD vectorization
        (\S\ref{sec:implementation}).
  \item Evaluation on three scientific benchmarks spanning neural networks,
        pharmacokinetics, and neuroimaging, demonstrating $2\times$ speedup
        over Julia with stronger static guarantees
        (\S\ref{sec:evaluation}).
  \item Integration with causal types for compile-time identifiability
        verification of interventional queries (\S\ref{sec:causal}).
\end{enumerate}

% ============================================================================
\section{Epistemic Type System}
\label{sec:type-system}
% ============================================================================

We formalize epistemic types as an extension of a simply-typed lambda calculus
with arithmetic.

\subsection{Syntax}

\paragraph{Types.}
\[
\begin{array}{rcl}
\tau & ::= & T \mid \knowfull{T}{\eps}{\Phi}{\mathit{tv}} \\
T & ::= & \texttt{i32} \mid \texttt{f64} \mid \texttt{bool} \mid \ldots \\
\eps & ::= & r \in \mathbb{R}_{\geq 0} \mid \alpha_\eps \quad \text{(uncertainty bound)} \\
\Phi & ::= & \texttt{Measured} \mid \texttt{Computed} \mid \texttt{Derived} \\
\mathit{tv} & ::= & r \in \mathbb{R}_{\geq 0} \mid \infty \quad \text{(temporal validity)}
\end{array}
\]

\paragraph{Terms.}
\[
\begin{array}{rcl}
e & ::= & x \mid c \mid e_1 \oplus e_2 \mid \texttt{measure}(e, \eps, \Phi, \mathit{tv}) \\
  & \mid & \texttt{extract}(e) \mid \texttt{transform}(f, e_1, \ldots, e_n)
\end{array}
\]
where $\oplus \in \{+, -, \times, \div\}$.

\subsection{Typing Rules}

The key innovation is that uncertainty propagation is encoded directly
in the typing rules, making GUM compliance a \emph{type system invariant}
rather than a runtime property.

\paragraph{Measurement introduction.}
\[
\inferrule*[right=T-Measure]
  {\Gamma \vdash e : T \\
   \eps \in \mathbb{R}_{\geq 0} \\
   \Phi \in \{\texttt{Measured}, \texttt{Computed}, \texttt{Derived}\}}
  {\Gamma \vdash \texttt{measure}(e, \eps, \Phi, \mathit{tv}) :
   \knowfull{T}{\eps}{\Phi}{\mathit{tv}}}
\]

\paragraph{Addition (GUM RSS propagation).}
For uncorrelated inputs, the GUM specifies root-sum-of-squares:
\[
\inferrule*[right=T-Add]
  {\Gamma \vdash e_1 : \knowfull{T}{\eps_1}{\Phi_1}{t_1} \\
   \Gamma \vdash e_2 : \knowfull{T}{\eps_2}{\Phi_2}{t_2} \\\\
   \eps_r = \sqrt{\eps_1^2 + \eps_2^2} \\
   t_r = \min(t_1, t_2)}
  {\Gamma \vdash e_1 + e_2 :
   \knowfull{T}{\eps_r}{\texttt{Computed}}{t_r}}
\]

\paragraph{Multiplication (relative uncertainty).}
\[
\inferrule*[right=T-Mul]
  {\Gamma \vdash e_1 : \knowfull{T}{\eps_1}{\Phi_1}{t_1} \\
   \Gamma \vdash e_2 : \knowfull{T}{\eps_2}{\Phi_2}{t_2} \\\\
   v_1 = \texttt{val}(e_1) \quad v_2 = \texttt{val}(e_2) \\
   \eps_r = |v_1 v_2| \sqrt{(\eps_1/v_1)^2 + (\eps_2/v_2)^2}}
  {\Gamma \vdash e_1 \times e_2 :
   \knowfull{T}{\eps_r}{\texttt{Computed}}{\min(t_1, t_2)}}
\]

\paragraph{General transformation (delta method).}
For $f : T_1 \times \cdots \times T_n \to U$:
\[
\inferrule*[right=T-Transform]
  {\Gamma \vdash e_i : \knowfull{T_i}{\eps_i}{\Phi_i}{t_i}
   \quad \forall i \in 1..n \\\\
   \eps_r = \sqrt{\textstyle\sum_{i=1}^{n}
     \left(\frac{\partial f}{\partial x_i}\right)^2 \eps_i^2}}
  {\Gamma \vdash \texttt{transform}(f, \bar{e}) :
   \knowfull{U}{\eps_r}{\texttt{Computed}}{\min_i t_i}}
\]

\paragraph{Extraction with confidence proof.}
\[
\inferrule*[right=T-Extract]
  {\Gamma \vdash e : \knowfull{T}{\eps}{\Phi}{\mathit{tv}} \\
   \texttt{SMT-PROVE}(\texttt{conf}(\eps) \geq \theta)}
  {\Gamma \vdash \texttt{extract}(e) : T}
\]
If the SMT solver cannot discharge the obligation, the compiler falls
back to a runtime check (gradual typing), ensuring that extraction is
always safe but potentially guarded.

\subsection{Operational Semantics}

Values are either scalar constants or knowledge tuples
$\langle v_0, \eps, \Phi, \mathit{tv} \rangle$.
The reduction rules mirror the typing rules:

\[
\inferrule*[right=E-Add]
  { }
  {\langle v_1, \eps_1, \Phi_1, t_1 \rangle + \langle v_2, \eps_2, \Phi_2, t_2 \rangle
   \longrightarrow
   \langle v_1{+}v_2, \sqrt{\eps_1^2{+}\eps_2^2}, \texttt{Computed}, \min(t_1,t_2) \rangle}
\]

\[
\inferrule*[right=E-Mul]
  {v_1, v_2 \neq 0}
  {\langle v_1, \eps_1, \Phi_1, t_1 \rangle \times \langle v_2, \eps_2, \Phi_2, t_2 \rangle
   \longrightarrow \\
   \langle v_1 v_2,\;
   |v_1 v_2| \sqrt{(\eps_1/v_1)^2{+}(\eps_2/v_2)^2},\;
   \texttt{Computed},\;
   \min(t_1,t_2) \rangle}
\]

\[
\inferrule*[right=E-Extract]
  {\texttt{conf}(\eps) \geq \theta}
  {\texttt{extract}(\langle v_0, \eps, \Phi, \mathit{tv} \rangle)
   \longrightarrow v_0}
\]

\subsection{Metatheory}

\begin{theorem}[Progress]
If $\emptyset \vdash e : \tau$, then either $e$ is a value or
$\exists e'.\ e \longrightarrow e'$.
\end{theorem}

\begin{proof}
By induction on the typing derivation.
The key case is T-Extract: if $e = \texttt{extract}(e')$
and $e'$ is a knowledge value with $\texttt{conf}(\eps) \geq \theta$,
then E-Extract applies.
If the confidence check fails, the gradual typing fallback
inserts a runtime guard, ensuring progress.
\end{proof}

\begin{theorem}[Preservation]
If $\Gamma \vdash e : \tau$ and $e \longrightarrow e'$,
then $\Gamma \vdash e' : \tau$.
\end{theorem}

\begin{proof}
By induction on the typing derivation and case analysis on
the reduction.
For E-Add: if $e_1 + e_2$ has type $\knowfull{T}{\eps_r}{\Phi'}{t'}$
with $\eps_r = \sqrt{\eps_1^2 + \eps_2^2}$, the result
$\langle v_1{+}v_2, \sqrt{\eps_1^2{+}\eps_2^2}, \texttt{Computed}, \min(t_1,t_2) \rangle$
has exactly this type.
\end{proof}

\begin{theorem}[GUM Compliance Soundness]
\label{thm:gum}
For any well-typed expression
$\Gamma \vdash e : \knowfull{T}{\eps}{\Phi}{\mathit{tv}}$
composed of arithmetic operations on epistemic values with uncorrelated
inputs, the computed uncertainty $\eps$ equals the GUM combined standard
uncertainty $u_c$ as defined by Equation~\eqref{eq:gum}.
\end{theorem}

\begin{proof}
By structural induction on $e$.
The base case (T-Measure) is immediate.
For T-Add: by IH, $\eps_i = u_c(e_i)$.
The GUM gives $u_c^2(e_1{+}e_2) = (\partial f/\partial x_1)^2 u^2(e_1)
+ (\partial f/\partial x_2)^2 u^2(e_2) = u^2(e_1) + u^2(e_2)$
since $\partial(x_1{+}x_2)/\partial x_i = 1$.
The typing rule gives $\eps_r = \sqrt{\eps_1^2 + \eps_2^2}
= \sqrt{u_c^2(e_1) + u_c^2(e_2)} = u_c(e_1{+}e_2)$.
Multiplication and division follow analogously via relative uncertainty.
\end{proof}

\subsection{Extension: ODE Integration}

For ODE solvers, numerical uncertainty grows with integration time.
We add a temporal degradation rule:
\[
\inferrule*[right=T-ODE-Step]
  {\Gamma \vdash e : \knowfull{T}{\eps(t)}{\Phi}{\mathit{tv}} \\
   \eps(t{+}\Delta t) = \eps(t) \cdot (1 + \alpha \cdot \Delta t)}
  {\Gamma \vdash \texttt{ode\_step}(e, \Delta t) :
   \knowfull{T}{\eps(t{+}\Delta t)}{\texttt{Computed}}{\mathit{tv}{-}\Delta t}}
\]
where $\alpha$ is the solver-dependent stability constant
($\alpha \approx 1$ for forward Euler, $\alpha \ll 1$ for RK4).

% ============================================================================
\section{Implementation}
\label{sec:implementation}
% ============================================================================

We implement epistemic types in \sounio, a systems programming language
with algebraic effects, linear types, and dimensional analysis.

\subsection{Compiler Architecture}

The \sounio compiler is implemented in approximately 45{,}000 lines of Rust,
organized as a multi-stage pipeline:
\[
\text{Source} \to \text{Lexer} \to \text{Parser} \to \text{AST}
\xrightarrow{\text{Check}} \text{HIR}
\xrightarrow{\text{Lower}} \text{SIR}
\to \text{HLIR} \to \text{Backend}
\]

The key innovation is the \textbf{Scientific IR (SIR)}, an intermediate
representation that treats epistemic operations as first-class instructions
rather than library calls:

\begin{lstlisting}[caption={SIR epistemic instructions}]
enum SirInst {
  PropagateAdd { lhs: Id, rhs: Id, result: Id },
  PropagateMul { lhs: Id, rhs: Id, result: Id },
  PropagateDiv { lhs: Id, rhs: Id, result: Id },
  VarianceFuse { inputs: Vec<Id>, result: Id },
}
\end{lstlisting}

\subsection{SIR Optimizations}

Three optimization passes reduce the overhead of epistemic tracking:

\paragraph{Variance fusion.}
Consecutive epistemic additions are fused into a single RSS computation:
$u^2(x_1 + x_2 + x_3) = u^2(x_1) + u^2(x_2) + u^2(x_3)$,
eliminating intermediate square root and squaring operations.

\paragraph{Exact elision.}
Operations involving $\texttt{Knowledge::exact}(v)$ (zero uncertainty)
skip variance computation entirely, reducing to scalar arithmetic.

\paragraph{SIMD vectorization.}
Value and variance are computed in parallel using 2-wide SIMD
(AVX2 on x86-64, NEON on ARM), exploiting the fact that
$\know{\texttt{f64}}$ stores value and variance in adjacent memory.

\subsection{Memory Layout}

$\know{\texttt{f64}}$ occupies 40 bytes, stack-allocated:

\begin{center}
\begin{tabular}{lr}
\toprule
Field & Size \\
\midrule
\texttt{value} & 8 bytes (\texttt{f64}) \\
\texttt{variance} & 8 bytes (\texttt{f64}) \\
\texttt{confidence} & 16 bytes (Beta parameters) \\
\texttt{provenance} & 8 bytes (pointer) \\
\bottomrule
\end{tabular}
\end{center}

\noindent
No heap allocation is required for the core uncertainty pathway.
Provenance chains are lazily allocated only when regulatory
auditing is enabled.

\subsection{SMT Integration}

For dependent epistemic types with refinement predicates
(e.g., $\know{T}$ \texttt{where confidence} $\geq$ \texttt{0.95}),
the compiler invokes Z3 to discharge proof obligations at compile time.
When Z3 cannot decide within a configurable timeout, the compiler
falls back to gradual typing with a runtime assertion.

% ============================================================================
\section{Evaluation}
\label{sec:evaluation}
% ============================================================================

We evaluate \sounio on three axes:
(1)~correctness of uncertainty propagation,
(2)~performance relative to existing UQ libraries, and
(3)~expressiveness on real scientific computing problems.

\subsection{Micro-Benchmarks: Performance}

Table~\ref{tab:perf} compares \sounio against Python's
\texttt{uncertainties}~\cite{lebigot2024uncertainties} and Julia's
\texttt{Measurements.jl}~\cite{giordano2016measurements} on five
benchmarks (median of 10 runs on AMD Ryzen 9 7950X, 64 GB RAM).

\begin{table}[t]
\centering
\caption{Execution time (ms) for uncertainty propagation benchmarks.}
\label{tab:perf}
\begin{tabular}{lrrrr}
\toprule
Benchmark & Python & Julia & \sounio & Speedup \\
\midrule
GUM Chain (100K)       & 423   & 45   & \textbf{22}  & $1.9\times$ Julia \\
Monte Carlo (100K)     & 1{,}266 & 58 & \textbf{31}  & $1.9\times$ Julia \\
Cockcroft-Gault (10K)  & 133   & 12   & \textbf{6}   & $2.0\times$ Julia \\
Matrix (50$\times$50)  & 12    & 2.1  & \textbf{1.0} & $2.1\times$ Julia \\
Transcendental (10K)   & 111   & 8.5  & \textbf{4.2} & $2.0\times$ Julia \\
\bottomrule
\end{tabular}
\end{table}

\sounio achieves consistent $\sim$2$\times$ speedup over Julia, attributable
to: (1)~no wrapper object allocation (stack-allocated variance),
(2)~SIMD parallelization of value/variance computation, and
(3)~compile-time resolution of propagation rules.
The 10--40$\times$ speedup over Python reflects the fundamental
overhead of dynamic dispatch and heap allocation for wrapper types.

\paragraph{Compile-time overhead.}
Epistemic type checking adds 12\% to compilation time on average,
dominated by SIR optimization passes.
For non-epistemic code paths, overhead is zero since
epistemic passes are not invoked.

\subsection{ODE Solver Validation}
\label{sec:ode-validation}

To validate numerical fidelity, we cross-validate \sounio's forward Euler
solver against the analytical Bateman equation for two-compartment
pharmacokinetics~\cite{gibaldi1982pharmacokinetics}:
\begin{align}
  C_{\text{central}}(t) = \frac{D \cdot k_a}{k_a - k_e}
    \left(e^{-k_e t} - e^{-k_a t}\right)
\end{align}

\begin{table}[t]
\centering
\caption{ODE solver accuracy vs.\ analytical Bateman equation at $t = 24$\,h.}
\label{tab:ode}
\begin{tabular}{llrrr}
\toprule
Method & $\Delta t$ & $C_{\text{central}}$ & Analytical & Error \\
\midrule
Euler (240 steps)  & 0.1\,h  & 49.79\,mg & 50.40\,mg & 1.20\% \\
Euler (2400 steps) & 0.01\,h & 50.34\,mg & 50.40\,mg & 0.12\% \\
\bottomrule
\end{tabular}
\end{table}

The 10$\times$ error reduction from 10$\times$ smaller $\Delta t$
confirms first-order convergence, validating the compiler's numerical
fidelity across multi-thousand-iteration loops with mutable struct state.

\subsection{Benchmark 1: QNN-MNIST with Epistemic Uncertainty}

Quaternionic neural networks (QNNs)~\cite{parcollet2019quaternion} use
Hamilton products for richer feature interactions than real-valued networks.
We implement a QNN with epistemic uncertainty tracking through every
weight and activation, demonstrating that \sounio can propagate
uncertainty through a complete training loop.

\paragraph{Architecture.}
A 2-layer QNN: 2 quaternion inputs $\to$ 2 hidden (ReLU) $\to$ 3 output
classes, totaling 15 epistemic quaternion parameters (60 scalar weights,
each carrying GUM uncertainty).

\paragraph{Epistemic tracking.}
Weights are initialized as $\know{\texttt{Quaternion}}$ values with
initial uncertainty $u = 0.1$.
Each training step applies Bayesian uncertainty concentration:
$u_{\text{new}} = u / \sqrt{1 + 0.1 \cdot n_{\text{seen}}}$,
modeling the intuition that observing more data reduces epistemic
uncertainty about the true weight values.

\paragraph{GUM propagation through Hamilton product.}
The quaternion product $q_1 q_2$ requires 16 multiplications and 12
additions.
Uncertainty propagates through each operation via the T-Mul and T-Add
rules, yielding output predictions with non-zero uncertainty that
reflects accumulated weight uncertainty.

\paragraph{Results.}
After 5 epochs of training on 9 samples (3 per class):
\begin{itemize}
  \item Weight uncertainty decreases from $0.1$ to $7.3 \times 10^{-14}$
        (Bayesian concentration effect)
  \item Prediction uncertainty is non-zero for all test inputs
        (GUM propagation active through 28+ operations per forward pass)
  \item Softmax outputs sum to 1.0 within floating-point tolerance
  \item Predictive entropy $H = 1.098$ nats (near maximum for 3 classes),
        correctly reflecting high uncertainty on this toy problem
\end{itemize}

\noindent
This benchmark validates that \sounio's type system correctly propagates
uncertainty through complex algebraic operations (Hamilton products)
and that the GUM compliance theorem (Theorem~\ref{thm:gum}) holds
empirically.

\subsection{Benchmark 2: PBPK Causal Intervention}

Physiologically-based pharmacokinetic (PBPK) models predict drug
concentrations in the body.
We implement a 2-compartment model with causal graph encoding,
demonstrating the integration of epistemic and causal types.

\paragraph{Model.}
Oral drug absorption ($k_a = 1.0\,\text{h}^{-1}$,
$k_e = 0.1\,\text{h}^{-1}$, $V_d = 50\,\text{L}$)
simulated via forward Euler over 24 hours (100 time steps).

\paragraph{Causal graph.}
The pharmacokinetic mechanism is encoded as a causal DAG:
\[
\text{Dose} \to \text{Plasma} \to \text{Effect} \qquad
\text{Genotype} \to \text{Metabolism} \to \text{Plasma}
\]
The causal query $P(\text{Effect} \mid \text{do}(\text{Dose} = 100\,\text{mg}))$
is identifiable via the backdoor criterion with adjustment set
$\{\text{Genotype}, \text{Metabolism}\}$.

\paragraph{Epistemic decomposition.}
Uncertainty in the causal average treatment effect (ATE) is decomposed into
three components:
\begin{itemize}
  \item $\eps_{\text{stat}} = 1.96/\sqrt{n} \approx 0.062$ (statistical, $n = 1000$)
  \item $\eps_{\text{causal}} = 0.03$ (positivity assumption violation for rare genotypes)
  \item $\eps_{\text{model}} = 0.05$ (PBPK structural uncertainty)
  \item $\eps_{\text{total}} = \sqrt{\eps_{\text{stat}}^2 + \eps_{\text{causal}}^2
        + \eps_{\text{model}}^2} \approx 0.085$
\end{itemize}

\paragraph{Results.}
\begin{itemize}
  \item ATE is positive (higher dose $\to$ stronger effect)
  \item Genotype-specific effects are biologically plausible:
        poor metabolizers show stronger effect than extensive metabolizers,
        who show stronger effect than ultra-rapid metabolizers
        (slower clearance $\to$ higher drug exposure)
  \item PK parameters ($C_{\max}$, AUC, $T_{\max}$) fall within
        realistic clinical ranges
  \item Non-identifiable queries are rejected at compile time
        by the causal type system (see \S\ref{sec:causal})
\end{itemize}

\subsection{Benchmark 3: fMRI Causal Connectivity}

Functional magnetic resonance imaging (fMRI) connectivity
analysis estimates directed interactions between brain
regions~\cite{eklund2016cluster}.
We implement causal connectivity analysis for the visual cortex.

\paragraph{Model.}
A 7-region visual cortex DAG:
V1 $\to$ V2 $\to$ V4 $\to$ IT (ventral stream),
V1 $\to$ MT $\to$ MST (dorsal stream),
with a latent confounder V1 $\leftrightarrow$ V2
(shared thalamocortical input).

\paragraph{Identifiability analysis.}
The causal type system verifies:
\begin{itemize}
  \item V2 $\to$ V4: \textbf{identifiable} (backdoor set $\{$V1$\}$)
  \item V4 $\to$ IT: \textbf{identifiable} (backdoor set $\{$V2$\}$)
  \item V1 $\to$ V2: \textbf{not identifiable} (latent confounder)
  \item V1 $\to$ MT, MT $\to$ MST: \textbf{identifiable} (unconfounded)
\end{itemize}

\paragraph{Bootstrap uncertainty.}
50-resample bootstrap provides epistemic uncertainty on connectivity
estimates, with 95\% Fisher z-transformed confidence intervals.

\paragraph{Results.}
The system correctly rejects the non-identifiable V1 $\to$ V2 query
at compile time, while successfully estimating all identifiable
connections with bounded uncertainty.
This demonstrates that \sounio prevents a class of
scientific errors (claiming causal connectivity from confounded
observations) that traditional analysis tools cannot detect.

% ============================================================================
\section{Causal Type Integration}
\label{sec:causal}
% ============================================================================

Epistemic types compose with causal types to enable
\emph{causal-epistemic reasoning}: interventional queries that
carry uncertainty bounds verified at compile time.

\paragraph{Causal-epistemic type.}
\[
\cknow{T}{\eps}{\text{do}(X{=}x)}
\]
represents a causal effect estimate with epistemic uncertainty $\eps$,
identified via the do-calculus~\cite{pearl2009causality}.

\paragraph{Identifiability verification.}
The do-operator only type-checks when the causal effect is
identifiable from observational data:
\[
\inferrule*[right=T-Do]
  {\Gamma \vdash G : \texttt{CausalGraph} \\
   \texttt{SMT-PROVE}(\texttt{identifiable}(G, X, Y))}
  {\Gamma \vdash \texttt{do}(G, X{=}x) : \cknow{Y}{\eps_{\text{causal}}}{\text{do}(X)}}
\]

This is, to our knowledge, the first type system that integrates
causal identifiability with measurement uncertainty.

% ============================================================================
\section{Related Work}
\label{sec:related}
% ============================================================================

\paragraph{Uncertainty propagation libraries.}
Python's \texttt{uncertainties}~\cite{lebigot2024uncertainties} uses
automatic differentiation for GUM propagation via operator overloading.
Julia's \texttt{Measurements.jl}~\cite{giordano2016measurements} uses
dual numbers.
Both are runtime-only: a function accepting \texttt{Float64} silently
discards uncertainty.
\sounio makes this a type error.

\paragraph{Probabilistic programming.}
Stan~\cite{carpenter2017stan}, Pyro~\cite{bingham2019pyro}, and
Gen~\cite{cusumano2019gen} model full posterior distributions.
These are powerful but heavyweight for GUM-style propagation, which
requires only first-order Taylor expansion, not MCMC inference.
Epistemic types occupy a lighter-weight point in the design space,
suitable for engineering and metrology applications.

\paragraph{Dependent and refinement types.}
Idris~\cite{brady2013idris} and F*~\cite{swamy2016dependent} support
dependent types for proving correctness properties.
Liquid Haskell uses refinement types with SMT solvers.
These systems can express predicates on values but do not
natively model statistical uncertainty or GUM propagation.

\paragraph{Units of measure.}
F\#'s units of measure~\cite{kennedy2009units} provide compile-time
dimensional analysis.
\sounio combines units with epistemic uncertainty, allowing types
to express both physical dimensions and associated measurement error:
$\know{\texttt{kg}, \eps=0.01}$.

\paragraph{Automatic differentiation.}
JAX~\cite{bradbury2018jax} and PyTorch~\cite{paszke2019pytorch} compute
gradients via AD.
\sounio uses similar techniques for uncertainty propagation
(Equation~\eqref{eq:gum} requires partial derivatives),
but treats uncertainty as a type-level concern rather than a
numerical computation.

\paragraph{Interval arithmetic.}
Libraries like MPFI compute conservative bounds via interval
arithmetic.
Unlike intervals that bound worst-case errors, GUM models
\emph{statistical} uncertainty with probabilistic guarantees,
providing tighter estimates for real-world measurements.

% ============================================================================
\section{Conclusion}
\label{sec:conclusion}
% ============================================================================

We have presented epistemic types, a type system extension that makes
measurement uncertainty a first-class type-level concept.
Unlike existing library approaches, epistemic types provide
\emph{compile-time guarantees} that uncertainty is preserved through
computation, propagated correctly per the GUM standard, and
accompanied by regulatory-grade provenance.

Our implementation in \sounio demonstrates that these guarantees come
with minimal overhead: $2\times$ speedup over Julia's
\texttt{Measurements.jl} and 12\% compile-time increase.
Three scientific benchmarks---quaternionic neural networks,
pharmacokinetic modeling, and neuroimaging connectivity---validate
that epistemic types are expressive enough for real scientific
computing while catching entire classes of uncertainty-related errors
at compile time.

The integration with causal types enables a further novel guarantee:
compile-time verification that interventional queries are identifiable
from observational data, preventing spurious causal claims.

\paragraph{Future work.}
\begin{itemize}
  \item Mechanized proofs of metatheory in Coq or Lean
  \item Correlation tracking for dependent measurements
        (GUM Supplement 1~\cite{jcgm2008supplement1})
  \item GPU-accelerated epistemic computations
  \item Empirical calibration studies on production scientific codebases
\end{itemize}

\paragraph{Availability.}
\sounio is open source under the MIT license.

\bibliographystyle{ACM-Reference-Format}
\bibliography{../references}

\end{document}
